%% =============================================================
%%  chapters/07_fusion_architecture.tex
%% =============================================================
\chapter{Sensor Fusion Architecture and Training Pipeline}
\label{ch:fusion-architecture}

\section{Overall Architecture Design}
\label{sec:overall-architecture-design}
The proposed architecture follows a two-branch feature-level fusion
design, in which a camera branch and a pressure heatmap branch each
extract modality-specific feature representations that are combined by a
fusion module before being passed to four task-specific output heads for
occupancy, passenger classification, weight, and pose estimation. This
design directly implements the feature-level fusion strategy motivated in
Section~\ref{sec:fusion-levels} and the design decisions summarised in
Section~\ref{sec:design-decisions}.

\section{Camera Branch}
\label{sec:camera-branch}
The camera branch consists of a convolutional neural network backbone
that processes the preprocessed, ROI-cropped camera frame described in
Section~\ref{sec:camera-roi-extraction} and produces a compact spatial
feature embedding. A backbone pretrained on a large-scale image dataset
is used as initialisation to compensate for the comparatively limited
size of the in-cabin camera dataset, with the final classification layers
removed and replaced by a projection to the shared fusion feature
dimension.
%% TODO: Specify exact backbone architecture (e.g. ResNet, EfficientNet,
%% MobileNet variant) chosen for the embedded compute budget.

\section{Pressure Heatmap Branch}
\label{sec:pressure-branch}
The pressure heatmap branch applies a separate, smaller convolutional
network directly to the single-channel, spatially aligned pressure
heatmap produced by the preprocessing pipeline in
Section~\ref{sec:pressure-spatial-alignment}, producing a feature
embedding of the same dimensionality as the camera branch output. A
smaller network is used for this branch given the lower spatial
resolution and information density of pressure heatmaps compared to
camera images.
%% TODO: Specify exact pressure branch architecture (number of
%% convolutional layers, channel widths).

\section{Fusion Module}
\label{sec:fusion-module}

\subsection{Feature-Level Fusion Strategy}
\label{sec:feature-level-fusion-strategy}
The camera and pressure feature embeddings are combined using a learned
cross-modal attention fusion module, in which each modality's features
attend to the other before being concatenated and projected to a shared
fused representation, following the cross-modal attention design
demonstrated for camera-depth fusion in \cite{oliveira2022crossmodal}.
This fused representation is then passed to the multi-task output heads
described in Section~\ref{sec:multi-task-output-heads}.
%% TODO: Report exact attention mechanism dimensionality and number of
%% fusion layers used.

\subsection{Fusion Alternatives Considered and Rejected}
\label{sec:fusion-alternatives-rejected}
Simple feature concatenation followed by a fully connected layer was
implemented as a baseline fusion strategy but was found to underweight
the lower-dimensional pressure branch relative to the camera branch;
decision-level fusion, in which independent per-modality predictions are
averaged, was also considered but rejected because it cannot model
cross-modal correlations such as a camera-visible object combined with a
pressure signature inconsistent with a human occupant. Both alternatives
are retained as baselines for the ablation study reported in
Section~\ref{sec:ablation-study}.
%% TODO: Report quantitative comparison against these alternatives once
%% experiments are complete.

\section{Multi-Task Output Heads}
\label{sec:multi-task-output-heads}

\subsection{Occupancy Detection Head}
\label{sec:occupancy-head}
The occupancy detection head is a lightweight fully connected layer
followed by a sigmoid activation, producing a binary occupied/unoccupied
probability for the corresponding seating position.

\subsection{Passenger Classification Head}
\label{sec:classification-head}
The passenger classification head is a fully connected layer followed by
a softmax activation over the predefined passenger classes (e.g. adult,
child, child seat, object), sharing the fused feature representation with
the other task heads.

\subsection{Weight Estimation Head}
\label{sec:weight-head}
The weight estimation head is a small fully connected regression network
that outputs a continuous estimate of occupant weight in kilograms,
trained only on samples labelled as human occupants.

\subsection{Pose Estimation Head}
\label{sec:pose-head}
The pose estimation head outputs either a fixed set of body keypoint
coordinates or a discrete pose category, depending on the granularity of
pose annotation available in the dataset described in
Chapter~\ref{ch:data-acquisition}.
%% TODO: Finalise whether pose is treated as keypoint regression or
%% discrete pose classification, and update this subsection accordingly.

\section{Loss Function Design and Task Balancing}
\label{sec:loss-function-design}
Each task head is trained with a task-appropriate loss term: binary
cross-entropy for occupancy, categorical cross-entropy for passenger
classification, mean squared error for weight estimation, and a
keypoint or classification loss for pose estimation. The combined
multi-task training objective is a weighted sum of the individual task
losses,

\begin{equation}
\label{eq:multitask-loss}
\mathcal{L}_{\text{total}} = \lambda_{\text{occ}} \, \mathcal{L}_{\text{occ}}
+ \lambda_{\text{cls}} \, \mathcal{L}_{\text{cls}}
+ \lambda_{\text{wt}} \, \mathcal{L}_{\text{wt}}
+ \lambda_{\text{pose}} \, \mathcal{L}_{\text{pose}},
\end{equation}

\noindent where $\lambda_{\text{occ}}$, $\lambda_{\text{cls}}$,
$\lambda_{\text{wt}}$, and $\lambda_{\text{pose}}$ are task weighting
coefficients. These coefficients are tuned to balance the differing
scales and convergence rates of the four loss terms, as discussed in the
multi-task learning background in Section~\ref{sec:multi-task-learning}.
%% TODO: Report final values of the task weighting coefficients and
%% whether a fixed schedule or learned/uncertainty-based weighting scheme
%% was ultimately used.

\section{Training Configuration}
\label{sec:training-configuration}
The model is implemented in PyTorch and trained using the Adam optimiser
with a cosine learning-rate decay schedule, an initial learning rate
selected via a short learning-rate range test, and a batch size chosen to
fit the available GPU memory. Training is performed on a workstation
equipped with a CUDA-capable GPU, with early stopping based on validation
loss to reduce overfitting given the moderate dataset size described in
Section~\ref{sec:final-dataset-composition}.
%% TODO: Report exact learning rate, batch size, number of epochs, GPU
%% model, and total training time once experiments are finalised.

\section{Implementation Details}
\label{sec:implementation-details}
The camera and pressure branches, fusion module, and task heads are
implemented as a single end-to-end trainable PyTorch model, with
modality-specific data loaders that apply the preprocessing steps
described in Chapter~\ref{ch:preprocessing-sync} on the fly during
training. Model checkpoints corresponding to the best validation
performance are retained for the evaluation reported in
Chapter~\ref{ch:results}.
%% TODO: Note any relevant software versions, random seed handling, or
%% reproducibility measures taken.
